\documentclass[11pt]{article}
\usepackage[T1]{fontenc}
\usepackage{lmodern}
\usepackage[margin=1in]{geometry}
\usepackage{amsmath,amssymb,amsthm,mathtools}
\usepackage{graphicx}
\usepackage{xcolor}
\usepackage{microtype}
\usepackage[hidelinks]{hyperref}
\usepackage{enumitem}

\newtheorem{theorem}{Theorem}
\newtheorem{lemma}{Lemma}
\newtheorem{remark}{Remark}
\newcommand{\R}{\mathbb R}
\newcommand{\N}{\mathbb N}
\newcommand{\Lip}{\operatorname{Lip}}
\newcommand{\eps}{\varepsilon}

\title{An LUR Non-UR Ball Without the Uniformly Continuous\\
Quasiconvex Extension Property}
\author{Candidate partial result for arXiv:2603.05206\\
\small fa\_banach\_001, lane 14, agent\_lane\_14}
\date{17 August 2026}

\begin{document}
\maketitle

\begin{center}
\fbox{\parbox{0.91\linewidth}{\textbf{Status.}
Candidate partial result, likely valid, pending human review.  The theorem
below gives an explicit LUR but non-UR body for which the extension property
fails.  It does \emph{not} prove the source's universal necessity statement
for every LUR non-UR body.}}
\end{center}

\section{Source question and scope}

The source is C.~A. De Bernardi and L. Vesel\'y, \emph{Extending
quasiconvex functions from uniformly convex sets}, arXiv:2603.05206
(2026)~\cite{source}.  On page~2, after recording that uniform rotundity
(UR) is sufficient and failure of local uniform rotundity (LUR) is an
obstruction, the authors ask:

\begin{quote}
Is the UR property of a convex body $C$ necessary for every uniformly
continuous quasiconvex function on $C$ to extend to a uniformly continuous
quasiconvex function on the whole ambient space?  Equivalently, what happens
for bodies which are LUR but not UR?
\end{quote}

\begin{center}
\includegraphics[width=0.94\linewidth]{figures/open_problem_crop.png}
\end{center}

We show that the desired extension property already fails for a concrete
LUR non-UR unit ball.  Thus LUR alone is not sufficient.  What remains open
is whether \emph{every} non-UR body admits such a witness.

\section{Definitions}

A function $f$ on a convex set is \emph{quasiconvex} (QC) if each sublevel
$\{f\le r\}$ is convex.  A norm, or its closed unit ball, is \emph{locally
uniformly rotund} if, whenever $x,y_j$ are unit vectors and
$\|x+y_j\|\to2$, one has $\|x-y_j\|\to0$.  It is \emph{uniformly rotund}
if the latter implication is uniform in the base point: for every
$\delta>0$ there is $\eta>0$ such that unit vectors separated by at least
$\delta$ have midpoint norm at most $1-\eta$.

\section{Idea of the proof}

The proof of Theorem~3.2 in the source paper constructs a Lipschitz QC
function whose increasing strict sublevels are a fan of half-planes
converging to a flat boundary segment.  We put this same fan into a sequence
of increasingly square but strictly convex two-dimensional balls
$B_{\ell_{p_n}^2}$, where $p_n\to\infty$.  Their $\ell_2$-sum is LUR because
every individual block is LUR, but it is not UR because the blocks become
arbitrarily flat.

On the sum ball we take the supremum of the coordinate copies of the planar
witness.  The coordinate structure does two jobs: it keeps this supremum
quasiconvex, and it lets us recover the full planar witness by restriction to
one block.  Convexity then forces any extension's strict sublevel in that
block to remain on one side of the corresponding cutting line.  Across all
blocks this forces an order-one jump between pairs whose distances tend to
zero, contradicting uniform continuity.

\section{The planar input}

Let $D=[-1,1]^2$, $z=(0,1)$, and, for $k\ge1$, put
\[
  \eps_k=2^{-k},\qquad b_k=1-2^{-k},\qquad
  L_k(t)=1+\eps_k-\frac{\eps_k}{b_k}t,
\]
and $H_k=\{(t,s):s<L_k(t)\}$.  The following is the precise piece of the
source construction that we use.

\begin{lemma}[normalized planar fan]\label{lem:fan}
There are a Lipschitz QC function $q:D\to[0,1]$ and numbers
$0<\tau_2<\tau_3<\cdots<1$, with $\tau_k\to1$, such that
\[
 q(0,0)=q(z)=0
 \quad\text{and}\quad
 \{q<\tau_k\}\cap B_{\ell_p^2}=H_k\cap B_{\ell_p^2}
 \tag{4.1}\label{eq:fan}
\]
whenever $1<p<\infty$ and the line $s=L_k(t)$ crosses
$\operatorname{int}B_{\ell_p^2}$.
\end{lemma}

\begin{proof}
This is a normalized restriction of the construction in the proof of
\cite[Theorem~3.2]{source}.  After a linear rescaling of its flat boundary
segment, that proof uses the open half-planes $H_k$ inside
$A=\{s<1\}$ and supplies a Lipschitz QC $g$ on $\overline A$ and increasing
levels $\alpha_k$ satisfying
\[
  \{g<\alpha_k\}\cap A=H_k\cap A,\qquad
  g(0,0)<\alpha_1,\qquad g(z)\le\alpha_1.
\]
The construction comes from \cite[Lemma~2.9]{fanpaper}; the source proof
also notes $\alpha_1<\alpha_\infty:=\sup_k\alpha_k$.

Here $\alpha_\infty<\infty$.  Indeed, for every $k$ the point $t=1$ admits
an $s\in(L_k(1),1)$, hence a point of the compact square $D\cap A$ outside
$H_k$.  At that point $g\ge\alpha_k$, so
$\alpha_k\le\max_D g$.  The levels are strictly increasing because their
strict sublevels $H_k\cap A$ are distinct.  Define
\[
 q=\min\!\left\{1,\max\!\left\{0,
       \frac{g-\alpha_1}{\alpha_\infty-\alpha_1}\right\}\right\}\Big|_D,
 \qquad
 \tau_k=\frac{\alpha_k-\alpha_1}
                  {\alpha_\infty-\alpha_1}.
\]
A nondecreasing scalar composition preserves quasiconvexity, and the clamp
preserves Lipschitz continuity.  Thus $q$ has the asserted regularity and
low values.  Away from the top edge $s=1$, its strict $\tau_k$-sublevel is
exactly $H_k$.  A strictly convex $\ell_p^2$ ball meets that top edge only
at $z$; since $q(z)=0<\tau_k$ and $z\in H_k$, equality
\eqref{eq:fan} follows.  Finally $\tau_k\to1$ by definition of the
supremum.
\end{proof}

\section{Two elementary geometric lemmas}

\begin{lemma}[the cutting lines enter explicit blocks]\label{lem:inside}
For $n\ge1$, set $k=n+1$ and $p_n=2^{3k}$.  The line $s=L_k(t)$ crosses
$\operatorname{int}B_{\ell_{p_n}^2}$.
\end{lemma}

\begin{proof}
Write $\eps=2^{-k}\le1/4$ and take
\[
 t=\frac{b_k+1}{2}=1-\frac\eps2,\qquad
 s=L_k(t)=1-\frac{\eps^2}{2(1-\eps)}.
\]
Both coordinates lie strictly between zero and one.  Since $1-u\le e^{-u}$,
and $p_n=\eps^{-3}$,
\[
 t^{p_n}+s^{p_n}
 \le \exp\!\left(-\frac1{2\eps^2}\right)
   +\exp\!\left(-\frac1{2\eps(1-\eps)}\right)
 \le e^{-8}+e^{-8/3}<1.
\]
Thus $(t,s)$ is an interior point of the block ball lying on the line.
\end{proof}

\begin{lemma}[half-plane continuation]\label{lem:halfplane}
Let $K$ be a convex body in $\R^2$, let $H$ be an open half-plane whose
boundary line meets $\operatorname{int}K$, and let $U$ be open and convex.
If $U\cap K=H\cap K$, then $U\subset H$.
\end{lemma}

\begin{proof}
Write $H=\{\ell<0\}$ for an affine functional $\ell$, and choose
$x\in\{\ell=0\}\cap\operatorname{int}K$.  If $U$ contained a point
$y$ with $\ell(y)\ge0$, openness would give a nearby $y'\in U$ with
$\ell(y')>0$.  Choose $u\in H\cap\operatorname{int}K$ arbitrarily close
to $x$.  Then $u\in U$.  The segment $[u,y']$ lies in $U$, and, by taking
$u$ sufficiently close to $x$, a point of this segment just beyond the
line remains in $\operatorname{int}K$.  It belongs to $U\cap K$ but not
to $H\cap K$, a contradiction.
\end{proof}

\section{Main theorem}

\begin{theorem}\label{thm:main}
There exist a separable Banach space $X$, an LUR but non-UR closed unit ball
$C=B_X$, and a Lipschitz quasiconvex function $f:C\to\R$ which has no
uniformly continuous quasiconvex extension to $X$.
\end{theorem}

\begin{proof}
For $p_n=2^{3(n+1)}$ let $E_n=\ell_{p_n}^2$ and put
\[
 X=\left(\bigoplus_{n=1}^{\infty}E_n\right)_{\!2},
 \qquad
 \|(x_n)\|_X=\left(\sum_n\|x_n\|_{p_n}^2\right)^{1/2}.
\]
This space is separable.

\smallskip\noindent
\emph{The ball is LUR.}
Every finite-dimensional $\ell_p^2$, $1<p<\infty$, is LUR.  For
completeness, the $\ell_2$-sum preserves this property as follows.  Suppose
$x=(x_n)$ and $y^j=(y_n^j)$ are unit vectors and $\|x+y^j\|\to2$.  The
vectors of component norms
\[
 a=(\|x_n\|)\quad\text{and}\quad
 c^j=(\|y_n^j\|)
\]
are unit vectors of $\ell_2$, and
$\|x+y^j\|\le\|a+c^j\|_2\le2$.  Uniform convexity of $\ell_2$ gives
$c^j\to a$.  Moreover,
\[
 \sum_n\left[(a_n+c_n^j)^2-\|x_n+y_n^j\|^2\right]\longrightarrow0.
\]
For each fixed $n$ with $a_n>0$, the $n$th nonnegative summand tends to
zero; after normalizing in $E_n$, LUR gives $y_n^j\to x_n$.  If $a_n=0$,
the same conclusion follows from $c_n^j\to0$.  Coordinate convergence,
together with $c^j\to a$ and a finite-head/small-tail argument, yields
$\|y^j-x\|_X\to0$.

\smallskip\noindent
\emph{The ball is not UR.}
In the $n$th block consider the unit vectors
\[
 u_n=(0,1),\qquad
 v_n=\left(\frac12,(1-2^{-p_n})^{1/p_n}\right).
\]
Their distance is at least $1/2$.  The second coordinate of their midpoint
tends to one, so by convexity
\[
 1\ge\left\|\frac{u_n+v_n}{2}\right\|_{p_n}
 \ge \frac{1+(1-2^{-p_n})^{1/p_n}}2\longrightarrow1.
\]
Embedding these vectors into the $n$th coordinate of $X$ violates uniform
rotundity.

\smallskip\noindent
\emph{The function.}
Let $q$ be supplied by Lemma~\ref{lem:fan} and define, for $x=(x_n)\in C$,
\[
 f(x)=\sup_n q(x_n).
\]
The square contains every block ball.  If $q$ has Euclidean Lipschitz
constant $M$, then
\[
 |q(r)-q(s)|\le M\|r-s\|_2
 \le \sqrt2M\|r-s\|_{p_n}
\]
for all $n$.  Consequently
\[
 |f(x)-f(y)|\le\sqrt2M\sup_n\|x_n-y_n\|_{p_n}
 \le\sqrt2M\|x-y\|_X.
\]
Also $\{f\le r\}=C\cap\prod_n\{q\le r\}$ is convex, so $f$ is QC.
Because $q(0)=0$ and $q\ge0$, the restriction of $f$ to the $n$th
coordinate ball is exactly $q$.

\smallskip\noindent
\emph{No uniformly continuous QC extension.}
Assume that $F:X\to\R$ is such an extension.  In the $n$th coordinate
plane set
\[
 U_n=\{w\in E_n:F(w)<\tau_{n+1}\}.
\]
This is open and convex, and Lemma~\ref{lem:fan} gives
\[
 U_n\cap B_{E_n}=H_{n+1}\cap B_{E_n}.
\]
By Lemma~\ref{lem:inside}, the cutting line meets the interior of the block
ball.  Lemma~\ref{lem:halfplane} therefore yields $U_n\subset H_{n+1}$.

Now regard the following as points in the $n$th coordinate plane:
\[
 z_n=(0,1),\qquad a_n=(0,1+\eps_n).
\]
At $t=0$, the boundary of $H_{n+1}$ has height
$1+\eps_{n+1}<1+\eps_n$, so $a_n\notin H_{n+1}$ and hence
\[
 F(a_n)\ge\tau_{n+1}\longrightarrow1.
\]
On the other hand, $z_n\in C$ and
$F(z_n)=f(z_n)=q(z)=0$, while
\[
 \|a_n-z_n\|_X=\eps_n\longrightarrow0.
\]
This contradicts uniform continuity of $F$.
\end{proof}

\section{Verification and limitations}

The script \texttt{code/verify\_geometry.py} checks the first twenty
explicit blocks in log space: the point in Lemma~\ref{lem:inside} is
strictly inside the block ball, the analytic exponential upper bound is
below one, $a_n$ is beyond the next cutting line, and the non-UR chord has
fixed separation.  These computations are not used as proof.

Five focused upgrades were attempted: moving fans along arbitrary bad
chords; supremum amalgamation without coordinates; projection and quotient
reductions; an ultrapower reduction; and Baire/diagonal uniformization of
extension moduli.  The common obstruction is localization.  In an arbitrary
LUR non-UR body the bad chords move, and no operation was found that combines
their local witnesses while preserving both quasiconvexity and uniformly
low values at all comparison points.  The $\ell_2$ sum solves exactly this
problem by coordinate isolation.  Thus Theorem~\ref{thm:main} is a proper
partial result and the universal source question remains open.

A bounded local-registry and arXiv/web search on 17 August 2026 used the
exact source title and the phrases ``quasiconvex extension LUR uniformly
rotund,'' ``quasiconvex extension locally uniformly rotund,'' and
``uniformly continuous quasiconvex extension.''  It recovered the source,
the authors' earlier positive UR theorem~\cite{ucpaper}, and general
LUR/non-UR renorming literature, but no later answer or matching block-sum
construction.  This is not an exhaustive priority search; novelty confidence
is moderate.

\paragraph{Human-review recommendation.}
Check the normalized planar-fan extraction in Lemma~\ref{lem:fan}, especially
the behavior at the top-edge point $z$, and the use of open strict sublevels
in Lemma~\ref{lem:halfplane}.  The remaining block-sum argument is elementary
and self-contained.

\begin{thebibliography}{9}
\bibitem{source}
C.~A. De Bernardi and L. Vesel\'y,
\emph{Extending quasiconvex functions from uniformly convex sets},
arXiv:2603.05206 (2026).

\bibitem{fanpaper}
C.~A. De Bernardi and L. Vesel\'y,
\emph{Rotundity properties, and non-extendability of Lipschitz quasiconvex
functions}, J. Convex Anal. \textbf{30} (2023), 329--342.

\bibitem{ucpaper}
C.~A. De Bernardi and L. Vesel\'y,
\emph{On extension of uniformly continuous quasiconvex functions},
Proc. Amer. Math. Soc. \textbf{151} (2023), 1705--1716.
\end{thebibliography}

\end{document}
